% Section content template for individual Educative lessons
% This file contains the actual content for a specific section
% Generated from Educative API section content

% Section: Requirements for the Elevator System
% Section ID: 6642858409066496
% Section Slug: requirements-for-the-elevator-system
% Generated: 2025-12-11T14:04:34.664986

In this lesson, we outline the functional and operational requirements for the elevator system. Identifying and understanding requirements is essential to defining the system's scope and ensuring a robust, user-friendly design.
\\
We'll use the notational convention to identify each requirement with a unique label ``Rn'', where ``R'' is short for requirement and ``n'' is a natural number.

\subsection{Requirement collection}\label{T3-iWeexW0GlqHAyijyXI}
\\
For the elevator design problem, the requirements are defined below:

% \vspace{-1.0\baselineskip}
\noindent
\begin{minipage}[t]{0.48\textwidth}\vspace{0pt}
\begin{itemize}
\item
\textbf{R1:} The elevator car system shall support a configurable number of floors (up to 15) and a configurable number of elevator cars (up to 3).
\item
\textbf{R2:} Each elevator car shall be capable of serving every floor and exist in one of four states: moving up, moving down, maintenance, or idle.
\item
\textbf{R3:} Elevator doors shall only open when the car is idle and auto-close after a configurable timeout unless the ``Open'' button is actively held.
\end{itemize}
\end{minipage}
\hfill
\begin{minipage}[t]{0.48\textwidth}\vspace{0pt}
\textit{Column component type 'MxGraphWidget' not yet supported.}
\end{minipage}

% \vspace{-1.0\baselineskip}
\noindent
\begin{minipage}[t]{0.48\textwidth}\vspace{0pt}
\begin{itemize}
\item
\textbf{R4:} Each floor shall have a panel with ``Up''/``Down'' call buttons, indicator lights, and an external display of the car's current floor and direction.
\item
\textbf{R5:} Each elevator car shall have buttons, ground (0) to 15th floor, for every floor, ``Open,'' ``Close,'' and an emergency-stop button, plus an internal display showing the current floor, direction, and load status.
\item
\textbf{R6:} Pressing the emergency-stop button shall immediately halt the car, keep doors closed, and alert building security/operators.
\item
\textbf{R7:}~Each elevator car shall enforce a maximum load of 680 kg, inhibit motion if this limit is exceeded, and emit an audible and visual alarm when overloaded.
\end{itemize}
\end{minipage}
\hfill
\begin{minipage}[t]{0.48\textwidth}\vspace{0pt}
\textit{Column component type 'MxGraphWidget' not yet supported.}
\end{minipage}

% \vspace{-1.0\baselineskip}
\noindent
\begin{minipage}[t]{0.48\textwidth}\vspace{0pt}
\begin{itemize}
\item
\textbf{R8:} The central controller shall assign the most appropriate car to each floor-call request, aiming to minimize average wait time, and command that car to move accordingly.
\item
\textbf{R9:} The elevator car system shall support calls from multiple passengers, with each passenger able to go to the same or different floors in the same or opposite direction.
\item
\textbf{R10:} The system shall support a maintenance state for each car in which:

\begin{itemize}
\item
\textbf{R10a:} The car is removed from normal dispatch (no new requests assigned).
\item
\textbf{R10b:} All elevator doors remain closed, and panel inputs are ignored.
\item
\textbf{R10c:} All UIs display ``Maintenance'' for that car.
\item
\textbf{R10d:} Upon exiting maintenance, the car returns to idle state at its current floor before resuming service.
\end{itemize}
\end{itemize}
\end{minipage}
\hfill
\begin{minipage}[t]{0.48\textwidth}\vspace{0pt}
\textit{Column component type 'MxGraphWidget' not yet supported.}
\end{minipage}

% \vspace{-1.0\baselineskip}
\noindent
\begin{minipage}[t]{0.48\textwidth}\vspace{0pt}
\begin{itemize}
\item
\textbf{R11:} All internal and external displays and indicators shall update in real time to reflect each elevator car's current floor, direction, and operational state (idle, moving, maintenance, or emergency).
\item
\textbf{R12:} The system shall communicate all operational states, errors, and safety messages to users via audio/visual indicators and displays, ensuring passengers are always aware of the elevator's status.
\end{itemize}
\end{minipage}
\hfill
\begin{minipage}[t]{0.48\textwidth}\vspace{0pt}
\textit{Column component type 'MxGraphWidget' not yet supported.}
\end{minipage}

We've identified the problem's requirements. In the next lesson, we will define different use cases for the elevator control system.

% End of section content